
\subsubsection{2.1 Entropy-Based Uncertainty Quantification}

Token-level entropy measures uncertainty in next-token prediction by computing the entropy of the softmax distribution over vocabulary. High entropy indicates the model is uncertain about which token to generate.

Kadavath et al. (2022) demonstrated that LLMs exhibit calibration properties, with confidence scores correlating with correctness on question-answering tasks. Kuhn et al. (2023) introduced semantic entropy, which clusters semantically equivalent answers before computing entropy. Malinin and Gales (2018) provided theoretical foundations for distinguishing epistemic and aleatoric uncertainty in neural networks.

\subsubsection{2.2 Consistency-Based Detection}

An alternative approach measures hallucination through response consistency across multiple generations. Manakul et al. (2023) introduced SelfCheckGPT, which generates multiple responses and computes consistency via embedding similarity. They evaluated on WikiBio biography generation. Wang et al. (2023) explored self-consistency in chain-of-thought reasoning.

\subsubsection{2.3 The Gap}

The two paradigms have evolved independently on different benchmarks:

\begin{table}[htbp]
\centering
\resizebox{\columnwidth}{!}{%
\begin{tabular}{llll}
\toprule
Method & Benchmark & Model & Direct Comparison \\
\midrule
Semantic Entropy (Kuhn 2023) & NLG tasks & Various & No consistency baseline \\
SelfCheckGPT (Manakul 2023) & WikiBio & GPT-3 & No entropy baseline \\
P(True) (Kadavath 2022) & QA tasks & Proprietary & No consistency baseline \\
\bottomrule
\end{tabular}
}
\end{table}

No prior work has compared both methods on the same factuality benchmark using the same model with analysis of whether signals are redundant or orthogonal.

